\documentclass[../../../include/open-logic-section]{subfiles}

\begin{document}

\olfileid[id]{sth}{story}{extensionality}
\olsection{Ekstensionalitas}

Hal pertama yang perlu dikatakan ialah bahwa himpunan dibedakan satu sama
lain berdasarkan !!{element}s-nya. Lebih tepatnya:

\begin{axiom}[Ekstensionalitas]
Jika himpunan $A$ dan $B$ mempunyai !!{element}s yang sama, maka $A$
dan $B$ adalah himpunan yang sama.
\[
  \lforall[A][\lforall[B][(\lforall[x][(x \in A \liff x \in B)] \lif
  \eq[A][B])]]
\]
\end{axiom}

Kita mengasumsikan hal ini di sepanjang \olref[sfr][][]{part}. Namun, hal
ini patut ditegaskan kembali. Aksioma Ekstensionalitas mengungkapkan gagasan
dasar bahwa suatu himpunan ditentukan oleh !!{element}s-nya. (Dengan
demikian, himpunan dapat dikontraskan dengan \emph{konsep}: objek-objek yang
persis sama dapat termasuk dalam banyak konsep yang berbeda.)

Mengapa kita menerima prinsip ini? Dapat dikatakan secara masuk akal bahwa
menolak Ekstensionalitas berarti memutuskan untuk meninggalkan segala sesuatu
yang bahkan masih layak disebut \emph{teori himpunan}. Teori himpunan tidak
lebih dan tidak kurang daripada teori koleksi ekstensional.

Namun, tantangan sesungguhnya dalam \olref[sth][][]{part} ialah menetapkan
prinsip-prinsip yang memberi tahu kita \emph{himpunan mana yang ada}. Ternyata,
satu-satunya jawaban yang benar-benar ``jelas'' atas pertanyaan ini dapat
dibuktikan keliru.

\end{document}
